\documentclass[12pt, a4paper]{article}

\usepackage[utf8]{inputenc}
\usepackage[T2A]{fontenc}
\usepackage[russian]{babel}
\usepackage{amsmath, amssymb, amsthm}
\usepackage{mathtools}
\usepackage{geometry}
\usepackage{xcolor}
\usepackage{tcolorbox}
\usepackage{booktabs}
\usepackage{array}
\usepackage{enumitem}
\usepackage{hyperref}
\usepackage{fancyhdr}
\usepackage{parskip}

\tcbuselibrary{skins, breakable}

\geometry{left=2.3cm, right=2.3cm, top=2.4cm, bottom=2.2cm}

% ===== ЦВЕТА =====
\definecolor{defblue}{RGB}{13, 71, 161}
\definecolor{defbluebg}{RGB}{227, 242, 253}
\definecolor{thmgreen}{RGB}{27, 94, 32}
\definecolor{thmgreenbg}{RGB}{232, 245, 233}
\definecolor{propviolet}{RGB}{74, 20, 140}
\definecolor{propvioletbg}{RGB}{243, 229, 245}
\definecolor{remamber}{RGB}{180, 60, 0}
\definecolor{remamberbg}{RGB}{255, 243, 224}
\definecolor{exgraybg}{RGB}{245, 245, 245}
\definecolor{exgray}{RGB}{60, 60, 60}
\definecolor{corollteal}{RGB}{0, 77, 64}
\definecolor{corolltbg}{RGB}{224, 242, 241}
\definecolor{proofbg}{RGB}{250, 250, 250}
\definecolor{qcolor}{RGB}{130, 0, 30}
\definecolor{fishkacol}{RGB}{136, 14, 79}
\definecolor{fishkabg}{RGB}{252, 228, 236}
\definecolor{razdelcol}{RGB}{26, 35, 126}

% ===== БЛОКИ =====
\tcbset{
  mybase/.style={
    breakable, enhanced,
    fonttitle=\bfseries\small,
    left=8pt, right=8pt, top=5pt, bottom=5pt,
    arc=4pt, boxrule=1pt,
    parbox=false
  }
}

\newtcolorbox{defbox}[1]{
  mybase, colback=defbluebg, colframe=defblue,
  colbacktitle=defblue, coltitle=white,
  title={\small Определение\ifx&#1&\else{}: {#1}\fi}
}
\newtcolorbox{thmbox}[1]{
  mybase, colback=thmgreenbg, colframe=thmgreen,
  colbacktitle=thmgreen, coltitle=white,
  title={\small Теорема\ifx&#1&\else{}: {#1}\fi}
}
\newtcolorbox{propbox}[1]{
  mybase, colback=propvioletbg, colframe=propviolet,
  colbacktitle=propviolet, coltitle=white,
  title={\small #1}
}
\newtcolorbox{rembox}[1]{
  mybase, colback=remamberbg, colframe=remamber,
  colbacktitle=remamber, coltitle=white,
  title={\small Замечание\ifx&#1&\else{}: {#1}\fi}
}
\newtcolorbox{exbox}[1]{
  mybase, colback=exgraybg, colframe=exgray,
  colbacktitle=exgray, coltitle=white,
  title={\small Пример\ifx&#1&\else{}: {#1}\fi}
}
\newtcolorbox{corbox}[1]{
  mybase, colback=corolltbg, colframe=corollteal,
  colbacktitle=corollteal, coltitle=white,
  title={\small Следствие\ifx&#1&\else{}: {#1}\fi}
}
\newtcolorbox{proofbox}[1]{
  mybase, colback=proofbg, colframe=gray!55,
  colbacktitle=gray!55, coltitle=black,
  title={\small #1}
}
\newtcolorbox{fishkabox}{
  mybase, colback=fishkabg, colframe=fishkacol,
  colbacktitle=fishkacol, coltitle=white,
  title={\small Фишка на собесе}
}

% ===== ЗАГОЛОВКИ =====
\newcommand{\razdel}[2]{%
  \clearpage%
  \noindent\colorbox{razdelcol}{\parbox{\dimexpr\linewidth-2\fboxsep\relax}{%
    \color{white}\LARGE\bfseries\strut\ \ Раздел #1. #2%
  }}%
  \vspace{10pt}%
}
\newcommand{\bilet}[2]{%
  \vspace{10pt}%
  \noindent\colorbox{qcolor}{\parbox{\dimexpr\linewidth-2\fboxsep\relax}{%
    \color{white}\large\bfseries\strut\ Билет #1. #2%
  }}%
  \nopagebreak\vspace{5pt}\par\nopagebreak%
}
\newcommand{\qsub}[1]{%
  \vspace{5pt}%
  {\bfseries\color{defblue} #1}%
  \par\vspace{2pt}%
}

% ===== КОЛОНТИТУЛЫ =====
\pagestyle{fancy}
\fancyhf{}
\fancyhead[L]{\small\color{gray} Классический ML --- конспект по билетам}
\fancyhead[R]{\small\color{gray} Подготовка к собесам}
\fancyfoot[C]{\small\thepage}
\renewcommand{\headrulewidth}{0.4pt}
\setlength{\headheight}{14pt}

% ===== МАТ. МАКРОСЫ =====
\newcommand{\E}{\mathbb{E}}
\newcommand{\R}{\mathbb{R}}
\newcommand{\Var}{\operatorname{Var}}
\newcommand{\cov}{\operatorname{cov}}
\newcommand{\sign}{\operatorname{sign}}
\newcommand{\rank}{\operatorname{rank}}
\newcommand{\cond}{\operatorname{cond}}
\newcommand{\RSS}{\operatorname{RSS}}
\newcommand{\diag}{\operatorname{diag}}
\newcommand{\tr}{^{\mathsf{T}}}
\newcommand{\norm}[1]{\left\lVert #1 \right\rVert}
\newcommand{\wh}[1]{\widehat{#1}}

\hypersetup{colorlinks=true, linkcolor=defblue, urlcolor=defblue, hidelinks}
\setlist[enumerate]{leftmargin=*, topsep=2pt, itemsep=2pt}
\setlist[itemize]{leftmargin=*, topsep=2pt, itemsep=2pt}

\begin{document}

\begin{titlepage}
\centering
\vspace*{4cm}
{\Huge\bfseries Классический ML}\\[0.6em]
{\Large конспект по билетам для собеса}\\[2cm]
{\large 14 разделов, 99 билетов}\\[0.4em]
{\large\color{gray} источник: \texttt{ml\_tickets.md}}\\[3cm]
{\color{gray} обновляется по ходу подготовки; \\ каждый билет: определения, выводы, численные примеры, фишки на собесах}
\end{titlepage}

% =====================================================================
\razdel{1}{Линейная регрессия}
% =====================================================================

\bilet{1}{Формализация задачи машинного обучения}

Любая задача обучения с учителем формализуется через тройку компонент.

\begin{defbox}{Три компоненты задачи обучения}
\begin{itemize}
  \item \textbf{Объекты} $x \in X$ --- сущности, про которые предсказываем (клиент, картинка, дом). Сами по себе абстрактны; модель работает не с объектом, а с его числовым описанием.
  \item \textbf{Признаки} (features) $f_j\colon X \to \R$ --- функции, сопоставляющие объекту измеримую характеристику. Вектор признаков $f(x) = (f_1(x),\ldots,f_m(x)) \in \R^m$.
  \item \textbf{Целевая переменная} $y \in Y$ --- то, что предсказываем. Существует истинная неизвестная зависимость $y\colon X \to Y$, которую восстанавливаем по данным.
\end{itemize}
\end{defbox}

\qsub{Типы признаков}
Бинарный ($\{0,1\}$), числовой ($\R$), категориальный (номинальный, неупорядоченное конечное множество), порядковый (ординальный, упорядоченное: школа $<$ бакалавр $<$ магистр).

\qsub{Матрица <<объект--признак>>}
Выборка из $n$ объектов даёт матрицу $X$ размера $[n \times (m+1)]$:
\[
  X = \begin{pmatrix}
    1 & x_{11} & \cdots & x_{1m}\\
    1 & x_{21} & \cdots & x_{2m}\\
    \vdots & & & \vdots\\
    1 & x_{n1} & \cdots & x_{nm}
  \end{pmatrix}, \qquad y = (y_1,\ldots,y_n)\tr.
\]
Строка $i$ --- объект $x_i$; столбец $j$ --- значения признака $f_j$ по всем объектам. \textbf{Первый столбец из единиц} $f_0 \equiv 1$ --- фиктивный признак под свободный член $w_0$ (bias). Отсюда размер $(m{+}1)$, а не $m$: это уловка, чтобы модель $f = w_0 + w_1 x_1 + \dots$ записывалась одним скалярным произведением $w\tr x$ без отдельного слагаемого.

\begin{defbox}{Типы обучения (по наличию $y$)}
\begin{itemize}
  \item \textbf{С учителем} (supervised): $y_i$ известны для всех обучающих объектов. Учимся по парам $(x_i, y_i)$.
  \item \textbf{Без учителя} (unsupervised): $y$ нет вообще. Ищем внутреннюю структуру --- кластеры, плотность, снижение размерности. Нет <<правильного ответа>>.
  \item \textbf{Частичное} (semi-supervised): $y$ известны у малой части объектов. Разметка дорогая, сырых данных много.
\end{itemize}
\end{defbox}

\qsub{Типы задач с учителем (по виду $Y$)}
\textbf{Регрессия} ($Y = \R$, непрерывное число: цена, спрос) $\to$ Раздел 1. \textbf{Классификация} ($Y$ --- конечное неупорядоченное множество меток: спам/не спам) $\to$ Раздел 5. \textbf{Ранжирование} ($Y$ --- порядок объектов: поисковая выдача) $\to$ Билет 91.

Ключевое отличие регрессии от классификации --- \textbf{тип $Y$}: непрерывный vs дискретный неупорядоченный. Это определяет функцию потерь, метрики и семейство моделей.

\begin{fishkabox}
Semi-supervised работает, только если структура $p(x)$ коррелирует с $p(y\mid x)$. Два предположения (см.\ билет 67):
\textbf{cluster assumption} --- объекты в одном плотном кластере имеют один $y$, граница классов проходит через разреженные области;
\textbf{manifold assumption} --- данные лежат на низкоразмерном многообразии, $y$ гладко меняется вдоль него.
Если $y$ не связан со структурой $p(x)$ (например, $y$ --- чётность случайного ID), неразмеченные данные не помогут вообще --- кластеры есть, но с $y$ не связаны.
\end{fishkabox}

\bilet{2}{Линейная регрессия: модель и метод наименьших квадратов}

\begin{defbox}{Модель линейной регрессии}
\[
  f(x, w) = w_0 + w_1 f_1(x) + \dots + w_m f_m(x) = w\tr x,
\]
где $x = (1, f_1(x), \ldots, f_m(x))\tr$, $w = (w_0, \ldots, w_m)\tr$. Для всей выборки: $\wh y = Xw$.
\end{defbox}

\begin{defbox}{Критерий МНК (метод наименьших квадратов)}
\[
  w = \arg\min_w \norm{y - Xw}_2^2 = \arg\min_w \sum_{i=1}^n (y_i - w\tr x_i)^2.
\]
\end{defbox}

Квадрат ошибки выбран не случайно: (1) даёт гладкую всюду дифференцируемую функцию (в отличие от $|\cdot|$); (2) сильнее штрафует большие отклонения (ошибка в 2 раза больше $\to$ вклад в 4 раза больше); (3) даёт аналитическое решение (билет 3) и статистическую интерпретацию как MLE при гауссовом шуме (билет 10).

\begin{propbox}{Линейность по весам, а не по $x$}
Модель <<линейна>>, потому что линейна по параметрам $w$, а НЕ по признакам $x$. Замена $x \to \varphi(x)$ (полиномы $\varphi = (x, x^2, x^3)$, RBF $\varphi = e^{-\norm{x-c}^2}$) оставляет модель линейной по $w$ (тот же МНК, та же машинерия), но как функция от исходного $x$ она нелинейна. \textbf{Полиномиальная регрессия} --- это линейная регрессия над признаками $\{x, x^2, \ldots, x^k\}$.
\end{propbox}

\begin{propbox}{Геометрическая интерпретация}
$y$ --- вектор в $\R^n$. Столбцы $X$ (каждый --- один признак по всем объектам) натягивают линейное подпространство (линейную оболочку) в $\R^n$. Модель $Xw$ для любого $w$ --- точка этого подпространства. МНК ищет \textbf{ближайшую точку подпространства к $y$} --- ортогональную проекцию $y$ на линейную оболочку столбцов $X$. Отсюда условие ортогональности невязки (билет 3).
\end{propbox}

\bilet{3}{Нормальное уравнение и матрица Грама}

\begin{proofbox}{Вывод нормального уравнения}
$\RSS(w) = \norm{y - Xw}^2 = (y - Xw)\tr(y - Xw) = y\tr y - 2 w\tr X\tr y + w\tr X\tr X w$. Это квадратичная функция с единственным минимумом; берём градиент и приравниваем к нулю:
\[
  \nabla \RSS(w) = -2 X\tr y + 2 X\tr X w = 0 \;\Longrightarrow\; \boxed{X\tr X\, w = X\tr y} \;\Longrightarrow\; w = (X\tr X)^{-1} X\tr y.
\]
\end{proofbox}

\begin{defbox}{Матрица Грама}
$X\tr X$ размера $[(m{+}1)\times(m{+}1)]$. Свойства:
\begin{itemize}
  \item \textbf{Симметрична}: $(X\tr X)\tr = X\tr X$.
  \item \textbf{Положительно полуопределена}: $\forall v\colon v\tr(X\tr X)v = (Xv)\tr(Xv) = \norm{Xv}^2 \ge 0$. Отсюда все собственные числа $\ge 0$ (пригодится в SVD, билет 6).
\end{itemize}
Матрица $(X\tr X)^{-1}X\tr$ --- псевдообратная Мура--Пенроуза (при полном ранге столбцов).
\end{defbox}

\begin{propbox}{Геометрический смысл нормального уравнения}
$X\tr(y - X\wh w) = 0$ покомпонентно означает: $x^{(j)\mathsf{T}}(y - \wh y) = 0$ для каждого столбца $x^{(j)}$, т.е.\ скалярное произведение $=0$ --- \textbf{невязка ортогональна каждому столбцу $X$}. Раз ортогональна всем векторам базиса плоскости --- ортогональна всей линейной оболочке. Это ровно определение ортогональной проекции: кратчайшее расстояние от точки $y$ до плоскости --- по перпендикуляру.
\end{propbox}

\qsub{Проблема обратимости и сложность}
$X\tr X$ обратима $\Leftrightarrow$ столбцы $X$ линейно независимы ($\rank X = m{+}1$). Ломается: (1) при \textbf{мультиколлинеарности} ($\rank X < m{+}1$); (2) при $n < m$ ($\rank X \le n < m{+}1$ по размерности). Обращение стоит $O(m^3)$ --- дорого при большом числе признаков; отсюда мотивация градиентного спуска (билет 5).

\begin{fishkabox}
Обратимость ломается по ДВУМ разным причинам: мультиколлинеарность (зависимость столбцов) и $n < m$ (объектов меньше признаков). Это не одно и то же: при $n \gg m$ с одним признаком-дубликатом матрица всё равно вырождена --- увеличение $n$ не лечит, т.к.\ проблема в ранге по СТОЛБЦАМ, а не по строкам.
\end{fishkabox}

\bilet{4}{Мультиколлинеарность и вырожденность матрицы Грама}

\begin{defbox}{Мультиколлинеарность}
Линейная зависимость (точная) или сильная корреляция (приближённая) между столбцами-признаками $X$.
\begin{itemize}
  \item \textbf{Точная} ($x_3 = 2x_1 + x_2$): $X\tr X$ строго вырождена, $\det = 0$, обратной нет.
  \item \textbf{Приближённая} (рост в см и в дюймах, $\rho \approx 0.999$): $X\tr X$ формально обратима, но плохо обусловлена ($\cond$ огромен).
\end{itemize}
\end{defbox}

\begin{propbox}{Взрыв весов и его смысл}
При точной зависимости проекция $\wh y = Xw$ НЕ меняется, но $w$ неединственно: бесконечно много комбинаций дают ту же $\wh y$ (можно забрать вес у $x_1$ и отдать эквивалент $x_3$). При приближённой --- решение единственно, но крайне чувствительно к шуму: малое возмущение $y$ или $X$ даёт огромный скачок $w$, у коррелирующих признаков появляются большие веса ПРОТИВОПОЛОЖНОГО знака (модель <<взаимно компенсирует>> два похожих признака). Это переобучение под шум конкретной выборки.
\end{propbox}

\begin{rembox}{$n<m$ --- отдельная причина}
Это НЕ мультиколлинеарность, а размерность: $\rank X \le \min(n, m{+}1)$. Если признаков больше объектов, ранг упирается в $n$, матрица Грама $(m{+}1)\times(m{+}1)$ гарантированно вырождена, даже если признаки не коррелируют (геномные данные: сотни образцов, десятки тысяч генов).
\end{rembox}

\qsub{Способы борьбы}
Отбор признаков; регуляризация (Ridge: $\tau E$ на диагональ гарантирует обратимость, сдвигая собственные числа на $+\tau$); PCA (декорреляция); псевдообратная через SVD (отбрасывание малых $\sigma$).

\begin{fishkabox}
\textbf{VIF} (Variance Inflation Factor): регрессируем $x_j$ на остальные признаки, $\text{VIF}_j = 1/(1 - R_j^2)$. VIF $> 5$--$10$ --- признак сильно объясняется остальными, кандидат на удаление.\\
Важно: мультиколлинеарность НЕ портит качество предсказаний ($\wh y$ та же), она портит интерпретируемость и устойчивость весов.
\end{fishkabox}

\bilet{5}{Градиентный спуск для линейной регрессии}

\begin{defbox}{Градиентный спуск}
\[
  w_{k+1} = w_k - h\,\nabla L(w_k), \qquad \nabla L(w) = 2 X\tr(Xw - y).
\]
$h$ --- learning rate. Итерация: считаем $\wh y = Xw_k$, невязку $(Xw_k - y)$, проецируем через $X\tr$ (вклад каждого признака в ошибку), шаг против градиента.
\end{defbox}

\qsub{Learning rate $h$}
$h$ большой $\to$ перепрыгивание минимума, расходимость. $h$ малый $\to$ медленная сходимость. Безопасная граница $h \le 1/L$, где $L = \lambda_{\max}(X\tr X)$ --- максимальная кривизна чаши потерь (билеты 17--18).

\qsub{Критерии останова}
$\norm{\nabla L(w_k)} < \varepsilon$ (почти на дне); $\norm{w_{k+1} - w_k} < \varepsilon$ (шаги микроскопические); лимит итераций.

\begin{exbox}{Пошаговый расчёт GD}
Данные: $x = (1,2,3)$, $y = (3,5,6)$, модель $w_0 + w_1 x$, старт $w = (0,0)$, $h = 0.1$.\\
\textbf{Шаг 1:} $\wh y = (0,0,0)$; невязка $r = \wh y - y = (-3,-5,-6)$;
$X\tr r = (-3{-}5{-}6,\; -3{-}10{-}18) = (-14, -31)$;
$\nabla L = 2(-14,-31) = (-28,-62)$;
$w \leftarrow (0,0) - 0.1(-28,-62) = (2.8,\, 6.2)$.\\
На сходимости $\nabla L = 2X\tr r = 0 \Leftrightarrow X\tr r = 0$ --- то же условие ортогональности, что нормальное уравнение достигает за один шаг, а GD --- постепенно.
\end{exbox}

\begin{table}[h]\centering\small
\begin{tabular}{@{}lll@{}}
\toprule
 & Нормальное уравнение & Градиентный спуск\\
\midrule
Точность & точное за один расчёт & приближённое, зависит от \# итераций\\
Сложность & $O(m^3)$ на обращение & $O(nm)$ за итерацию, итераций много\\
Обратимость $X\tr X$ & требуется & не требуется\\
Большие $n, m$ & плохо & хорошо (+ mini-batch, билет 19)\\
\bottomrule
\end{tabular}
\end{table}

\bilet{6}{Сингулярное разложение (SVD): определение и свойства}

\begin{defbox}{SVD}
Любую матрицу $X$ $[n\times m]$ можно разложить:
\[
  X = U S V\tr, \quad U\tr U = I\ (n\times n),\quad V\tr V = I\ (m\times m),\quad S = \diag(\sigma_1 \ge \dots \ge \sigma_r \ge 0),
\]
$U$ --- левые, $V$ --- правые сингулярные векторы, $\sigma_j$ --- сингулярные числа, $r = \rank X$.
\end{defbox}

\begin{rembox}{}
SVD существует для ЛЮБОЙ матрицы (прямоугольной, вырожденной). Собственное разложение --- только для квадратных диагонализуемых. В этом огромное практическое преимущество SVD.
\end{rembox}

\begin{proofbox}{Связь с матрицей Грама}
\[
  X\tr X = (USV\tr)\tr(USV\tr) = V S\tr U\tr U S V\tr = V(S\tr S)V\tr = V\diag(\sigma_j^2)V\tr.
\]
Использовали $U\tr U = I$. Это форма собственного разложения: $\sigma_j^2$ --- собственные числа $X\tr X$, столбцы $V$ --- её собственные векторы. Отсюда $\sigma_j^2 \ge 0$ (полуопределённость). Мостик к PCA (билет 31): там $(X\tr X)\text{Rot} = \lambda\,\text{Rot}$ даёт $\text{Rot} = V$, $\lambda = \sigma^2$.
\end{proofbox}

\begin{propbox}{Геометрия: поворот--растяжение--поворот}
Умножение на $X = USV\tr$ = три действия: (1) поворот $V\tr$ в пространстве признаков; (2) растяжение по осям на $\sigma_1, \ldots, \sigma_r$; (3) поворот $U$ в пространстве объектов. Любое линейное отображение = поворот $\to$ растяжение $\to$ поворот.
\end{propbox}

\qsub{Усечённое SVD и обусловленность}
$X_k = \sum_{j=1}^k \sigma_j u_j v_j\tr$ --- оставляем top-$k$, отбрасываем остальные (наилучшее приближение ранга $k$, теорема Эккарта--Янга, билет 30). Число обусловленности $\cond(X) = \sigma_{\max}/\sigma_{\min} = \sigma_1/\sigma_r$.

\begin{fishkabox}
При усечении (оставляем top-$k$): $\cond(X_k) = \sigma_1/\sigma_k \le \sigma_1/\sigma_r = \cond(X)$ --- обусловленность УМЕНЬШАЕТСЯ (устойчивость растёт). Не путай направление: чем БОЛЬШЕ $\cond$, тем ХУЖE устойчивость.
\end{fishkabox}

\bilet{7}{Решение МНК через SVD и псевдообратная матрица}

\begin{proofbox}{Вывод}
Подставим $X = USV\tr$ в $w = (X\tr X)^{-1}X\tr y$. Знаем $X\tr X = V(S\tr S)V\tr$, значит $(X\tr X)^{-1} = V(S\tr S)^{-1}V\tr$ (т.к.\ $V$ ортогональна). $X\tr = VS\tr U\tr$. Тогда
\[
  w = V(S\tr S)^{-1}V\tr \cdot V S\tr U\tr y = V(S\tr S)^{-1}S\tr U\tr y = \sum_{j=1}^r \frac{1}{\sigma_j} V_j (U_j\tr y).
\]
\end{proofbox}

\begin{defbox}{Псевдообратная Мура--Пенроуза}
$X^+ = V S^+ U\tr$, $S^+$ --- $S$ с ненулевыми $\sigma_j$, заменёнными на $1/\sigma_j$. $w = X^+ y$. Всегда определена, даже при вырожденной $X\tr X$. При полном ранге совпадает с $(X\tr X)^{-1}X\tr y$; при вырожденности даёт решение \textbf{минимальной нормы} $\norm{w}$ среди всех $w$ с одинаковой (минимальной) ошибкой.
\end{defbox}

\qsub{Смысл формулы по шагам}
$U_j\tr y$ --- проекция $y$ на $j$-е направление (сколько $y$ вдоль него). Делим на $\sigma_j$ --- обратное растяжение (компенсируем, что $X$ растягивал в $\sigma_j$ раз). Умножаем на $V_j$ --- переносим в пространство весов.

\begin{exbox}{Почему малое $\sigma$ опасно}
Пусть $\sigma_1 = 10$, $\sigma_2 = 0.001$, шум в $y$ даёт $U_j\tr y = 0.5$ на обоих направлениях. Вклад в $w$: $(1/10)\cdot 0.5 = 0.05$ (норм) против $(1/0.001)\cdot 0.5 = 500$ (взрыв!). Одинаковый шум усилен в 10000 раз из-за деления на почти-ноль. Усечение (отбрасывание $\sigma_2$) убирает этот взрыв.
\end{exbox}

\begin{fishkabox}
$\cond(X\tr X) = \cond(X)^2$, т.к.\ собственные числа $X\tr X$ --- это $\sigma_j^2$. Значит явное обращение $X\tr X$ численно вдвое хуже, чем работа напрямую с $X$ через SVD. При $\cond(X) = 100$ имеем $\cond(X\tr X) = 10000$. Поэтому \texttt{numpy.lstsq}, \texttt{sklearn} решают МНК через SVD/QR, а не через обращение матрицы Грама.
\end{fishkabox}

\bilet{8}{Метрики качества регрессии}

\begin{rembox}{Функция потерь $\ne$ метрика}
Функция потерь оптимизируется при обучении (MSE в МНК --- дифференцируема, даёт аналитику). Метрика измеряет результат постфактум и может быть другой (обучали через MSE, отчитываемся через MAE ради интерпретируемости).
\end{rembox}

\begin{defbox}{Метрики}
\begin{itemize}
  \item $\text{MSE} = \frac1n\sum(y_i - \wh y_i)^2$ --- квадратичный штраф, не робастна к выбросам, единицы $y^2$.
  \item $\text{RMSE} = \sqrt{\text{MSE}}$ --- единицы $y$ (интерпретируема), всё ещё чувствительна к выбросам.
  \item $\text{MAE} = \frac1n\sum|y_i - \wh y_i|$ --- линейный штраф, робастна, не диф.\ в 0. Всегда $\text{MAE} \le \text{RMSE}$.
  \item $\text{MAPE} = \frac1n\sum \frac{|y_i - \wh y_i|}{|y_i|}$ --- в \%, масштабно-независима, но ломается при $y_i \approx 0$ и НЕ робастна (чувствительна через знаменатель).
  \item $R^2 = 1 - \text{SS}_{\text{res}}/\text{SS}_{\text{tot}}$, $\text{SS}_{\text{res}} = \sum(y_i-\wh y_i)^2$, $\text{SS}_{\text{tot}} = \sum(y_i - \bar y)^2$ --- доля объяснённой дисперсии. $R^2 = 1$ идеал, $R^2 = 0$ = константный прогноз $\bar y$, $R^2 < 0$ хуже среднего.
\end{itemize}
\end{defbox}

\begin{propbox}{Диагностика по соотношению RMSE/MAE}
$\text{RMSE} \gg \text{MAE}$ --- есть отдельные большие выбросы (квадрат раздувает RMSE). $\text{RMSE} \approx \text{MAE}$ --- ошибки распределены равномерно. Для прода, где критичны редкие катастрофические ошибки, смотрят RMSE.
\end{propbox}

\begin{defbox}{Adjusted $R^2$}
\[
  R^2_{\text{adj}} = 1 - (1 - R^2)\frac{n-1}{n-p-1}, \quad p = \text{\# признаков}.
\]
\end{defbox}

\begin{fishkabox}
Обычный $R^2$ НИКОГДА не убывает при добавлении признаков (даже случайного шума) --- у модели больше степеней свободы, в худшем случае вес нового признака $= 0$. Adjusted $R^2$ штрафует за $p$: при $p \to n$ знаменатель $(n-p-1) \to 0$, дробь $\to \infty$, $R^2_{\text{adj}} \to -\infty$ --- жёстко сигнализирует о переобучении, даже если обычный $R^2 \approx 1$.
\end{fishkabox}

\bilet{9}{Предобработка признаков: масштабирование и кодирование}

\begin{defbox}{Масштабирование числовых признаков}
\textbf{Стандартизация} (z-score): $x' = (x - \mu)/\sigma$ --- центр 0, дисперсия 1. \textbf{Min-max}: $x' = (x - \min)/(\max - \min)$ --- диапазон $[0,1]$, чувствительна к выбросам.
\end{defbox}

\qsub{Зачем масштабировать}
(1) \textbf{Ускоряет GD}: разные масштабы признаков $\to$ разброс собственных чисел $X\tr X$ $\to$ плохая обусловленность $\to$ вытянутая эллиптическая чаша $\to$ зигзаг (билет 5/11). Масштабирование делает чашу круглее. (2) \textbf{Методы на расстояниях} (kNN, k-means, SVM-RBF): без масштабирования признак с большим диапазоном доминирует в евклидовом расстоянии. (3) \textbf{Сопоставимость весов} и справедливость регуляризации (иначе штраф сильнее давит на признаки малого масштаба).

\begin{defbox}{Кодирование категорий}
\textbf{One-hot}: столбец на значение. \textbf{Ordinal}: номер по порядку (ложный порядок для неупорядоченных). \textbf{Target/mean encoding}: категория $\to$ среднее $y$ по ней (риск утечки). \textbf{Hashing}: хэш в фикс.\ число столбцов (коллизии). \textbf{Embeddings}: обучаемый плотный вектор.
\end{defbox}

\begin{propbox}{Dummy trap}
Если оставить все $K$ one-hot столбцов: их сумма $= $ вектор единиц $= f_0$ (фиктивный признак) $\Rightarrow$ точная линейная зависимость $\Rightarrow$ мультиколлинеарность, $X\tr X$ вырождена. Решение: убрать один столбец ($K{-}1$ достаточно).
\end{propbox}

\begin{rembox}{Деревья}
Деревьям масштабирование НЕ нужно: пороговые сплиты $[f_j \gtrless d_j]$ инвариантны к монотонным преобразованиям. По той же причине ordinal/target encoding для деревьев ОК даже без <<настоящего>> порядка.
\end{rembox}

\begin{fishkabox}
Target encoding без out-of-fold --- утечка: для объекта его собственный $y_i$ участвует в его же признаке. Правильно: считать среднее по остальным фолдам + сглаживание (взвешивание с глобальным средним для редких категорий). Формально сглаживание = байесовская априорная оценка (см.\ билет 45, сглаживание Лапласа).
\end{fishkabox}

\bilet{10}{Вероятностный взгляд на линейную регрессию}

\begin{defbox}{Вероятностная модель}
$y = w\tr x + \varepsilon$, $\varepsilon \sim \mathcal{N}(0, \sigma^2)$, т.е.\ $p(y\mid x, w) = \mathcal{N}(y \mid w\tr x, \sigma^2)$: среднее линейно по $x$, дисперсия $\sigma^2$ постоянна (гомоскедастичность).
\end{defbox}

\begin{proofbox}{МНК = максимум правдоподобия при гауссовом шуме}
$L(w) = \prod_i \frac{1}{\sqrt{2\pi\sigma^2}} e^{-(y_i - w\tr x_i)^2/2\sigma^2}$. Логарифм:
\[
  \log L(w) = -\tfrac{n}{2}\log(2\pi\sigma^2) - \tfrac{1}{2\sigma^2}\sum_i (y_i - w\tr x_i)^2.
\]
Первое слагаемое не зависит от $w$, значит $\arg\max_w \log L = \arg\min_w \sum(y_i - w\tr x_i)^2$ --- это МНК.
\end{proofbox}

\begin{propbox}{Регуляризация как байесовская MAP-оценка}
$p(w \mid y, X) \propto p(y \mid X, w)\, p(w)$. Максимизация $\log$-апостериора $= $ минимизация $\frac{1}{2\sigma^2}\norm{Xw - y}^2 + [-\log p(w)]$:
\begin{itemize}
  \item \textbf{Ridge} $=$ гауссово априорное $p(w) = \mathcal{N}(0, \sigma_0^2 I)$: $-\log p(w) \propto \tau\norm{w}^2$, $\tau = \sigma^2/\sigma_0^2$.
  \item \textbf{Lasso} $=$ лапласово априорное $p(w) \propto e^{-\lambda\sum|w_j|}$: острый пик в 0 $\to$ зануление.
  \item \textbf{MAE} $=$ лапласовский шум (тяжёлые хвосты $\to$ робастность к выбросам).
\end{itemize}
$\tau$ = отношение дисперсии шума к дисперсии априорного: чем увереннее, что веса малы (малое $\sigma_0^2$), тем сильнее регуляризация.
\end{propbox}

\qsub{Предположения модели}
Независимость наблюдений; гауссов шум; гомоскедастичность ($\sigma^2$ не зависит от $x$).

\begin{fishkabox}
Гетероскедастичность (дисперсия шума растёт с $x$): МНК-оценка $w$ ОСТАЁТСЯ несмещённой (несмещённость требует лишь $\E[\varepsilon\mid X] = 0$, не постоянства $\sigma^2$). Ломается ЭФФЕКТИВНОСТЬ: по теореме Гаусса--Маркова МНК = BLUE только при гомоскедастичности; иначе есть оценки с меньшей дисперсией (WLS). Главное --- стандартные ошибки $\Var(w) = \sigma^2(X\tr X)^{-1}$ становятся неверными $\to$ доверительные интервалы и $p$-value некорректны. Лечение: робастные ошибки (White/Huber).
\end{fishkabox}

% =====================================================================
\razdel{2}{Регуляризация и теория оптимизации}
% =====================================================================

\bilet{11}{Число обусловленности и устойчивость решения}

\begin{defbox}{Число обусловленности}
\[
  \cond(A) = \norm{A}\cdot\norm{A^{-1}} = \frac{\sigma_{\max}}{\sigma_{\min}}, \qquad \frac{\norm{\Delta x}}{\norm{x}} \le \cond(A)\,\frac{\norm{\Delta b}}{\norm{b}}.
\]
$\cond$ --- коэффициент усиления относительной ошибки входа в решении $Ax = b$ (у нас $A = X\tr X$).
\end{defbox}

\begin{propbox}{Геометрия <<почти параллельных прямых>>}
Решение системы = пересечение прямых. Под прямым углом --- точка чёткая. Почти параллельны --- малейшее дрожание уносит пересечение далеко. Малый угол между строками $A$ $\Rightarrow$ почти линейная зависимость строк $\Rightarrow$ вырождается ПОСЛЕДНЕЕ независимое направление $\Rightarrow$ $\sigma_{\min} \to 0$ (не $\sigma_{\max}$: она определяется длиной строк). Для матрицы из единичных строк под углом $\theta$: $\cond \approx 2/\theta$ при малых $\theta$.
\end{propbox}

\begin{fishkabox}
Одна цепочка эквивалентностей: мультиколлинеарность $\Rightarrow$ $\sigma_{\min} \to 0$ $\Rightarrow$ $\cond$ огромен $\Rightarrow$ коэффициент $1/\sigma_{\min}$ в решении взрывается $\Rightarrow$ веса неустойчивы к шуму. Высокий $\cond$ $\Leftrightarrow$ мультиколлинеарность $\Leftrightarrow$ взрыв весов --- одно явление, три описания. Лечится Ridge ($\tau E$ сдвигает $\lambda_{\min} \to \lambda_{\min}+\tau$) и масштабированием.
\end{fishkabox}

\bilet{12}{Переобучение и недообучение}

\begin{defbox}{Определения через симптомы}
\textbf{Переобучение} (overfitting): подстройка под шум train; ошибка на train мала (иногда $\approx 0$), на test велика; симптом --- огромные неустойчивые веса (high variance). \textbf{Недообучение} (underfitting): модель слишком проста; ошибка велика и на train, и на test (high bias).
\end{defbox}

\qsub{Train / Validation / Test}
Train --- считаются веса. Validation --- подбираются гиперпараметры ($\tau$, $K$, глубина). Test --- используется ОДИН раз в конце для честной оценки. Подбор гиперпараметров по test превращает его во вторую validation $\to$ оптимистичное смещение (утечка, билет 23).

\begin{fishkabox}
\textbf{Early stopping = implicit регуляризация} без штрафного члена в loss. Механизм: норма весов $\norm{w}$ монотонно РАСТЁТ по мере итераций GD (старт с $w \approx 0$). Обрыв обучения не даёт $\norm{w}$ дорасти до величины несжатой МНК-оценки --- функционально имитирует ограничение на $\norm{w}$. Для линейной регрессии: раннее остановленный GD $\approx$ Ridge с $\tau \propto 1/(\text{\# итераций})$. Для бустинга: меньше деревьев = меньше эффективная сложность.
\end{fishkabox}

\bilet{13}{Ridge-регрессия (L2-регуляризация)}

\begin{defbox}{Ridge}
\[
  L = \tfrac1n\norm{Xw - y}^2 + \tau\norm{w}_2^2, \qquad w = (X\tr X + \tau E)^{-1}X\tr y.
\]
\end{defbox}

\begin{proofbox}{Механизм добавки $\tau E$}
Собственные числа $\lambda_i$ матрицы $X\tr X$ (все $\ge 0$) переходят в $\lambda_i + \tau > 0$ у $(X\tr X + \tau E)$. Матрица становится строго положительно определённой $\Rightarrow$ обратима всегда. $\cond = (\lambda_{\max}+\tau)/(\lambda_{\min}+\tau)$ падает --- лечится и обусловленность. \textbf{Ridge НЕ убирает корреляцию признаков} (данные не трогает), только сдвигает собственные числа --- не путать с декорреляцией через PCA.
\end{proofbox}

\qsub{Свойства}
$w_0$ (bias) обычно не регуляризуют (это сдвиг, не сложность). Ridge сжимает веса к нулю, но НЕ зануляет точно (обращение матрицы не даёт точный ноль при $X\tr y \ne 0$). Влияние $\tau$: рост $\to$ bias$\uparrow$, variance$\downarrow$.

\begin{fishkabox}
Чтобы по весам Ridge понять важность признаков --- сначала СТАНДАРТИЗИРОВАТЬ $X$ (иначе малый вес может значить <<признак большого масштаба>>, а не <<неважный>>). Ridge даёт непрерывное ранжирование, но не автоматический отбор --- для отбора нужен Lasso.
\end{fishkabox}

\bilet{14}{Lasso-регрессия (L1-регуляризация)}

\begin{defbox}{Lasso}
\[
  L = \tfrac1n\norm{Xw - y}^2 + \lambda\norm{w}_1, \quad \norm{w}_1 = \sum_j |w_j|.
\]
Главное свойство: ЗАНУЛЯЕТ веса неинформативных признаков $\to$ автоматический отбор признаков (sparsity), интерпретируемость.
\end{defbox}

\qsub{Нет аналитики}
$|w_j|$ не дифференцируема в $0$ (излом: слева $-1$, справа $+1$). Нужен субградиент (в точке излома --- весь отрезок $[-1,1]$). Решают координатным спуском.

\begin{propbox}{Soft-thresholding --- механизм зануления}
\[
  w_j^* = S_\lambda(\rho_j) = \sign(\rho_j)\cdot\max(|\rho_j| - \lambda,\ 0),
\]
$\rho_j$ --- <<сырой>> МНК-вклад признака $j$. Из $|\rho_j|$ вычитаем $\lambda$; если результат отрицателен --- обрезаем в ноль. Сравни с Ridge: там УМНОЖЕНИЕ на коэффициент $< 1$ (ноль недостижим), здесь ВЫЧИТАНИЕ-с-обрезкой $\to$ точный ноль при $|\rho_j| \le \lambda$.
\end{propbox}

\begin{exbox}{Soft-threshold}
$\rho_j = 0.3$, $\lambda = 0.5$: $\max(0.3 - 0.5, 0) = \max(-0.2, 0) = 0$. Сигнал слабее порога --- зануляется.
\end{exbox}

\begin{fishkabox}
Нестабильность Lasso при коллинеарной группе --- НЕ из-за порядка обхода координат (координатный спуск сходится к глобальному оптимуму независимо от порядка). Причина: при почти идентичных признаках задача имеет ВЫРОЖДЕННОЕ решение --- <<весь вес на A>> и <<весь вес на B>> дают почти одинаковый loss. Геометрически (билет 15): эллипсы RSS касаются ромба вдоль целого ребра, и в какой угол упрётся решение --- решает шум. У Ridge решение гладко и единственно $\Rightarrow$ вес делится между коррелирующими признаками стабильно.
\end{fishkabox}

\bilet{15}{Геометрия регуляризации: L1 vs L2}

\begin{propbox}{Эквивалентная условная задача}
$\min[\RSS(w) + \lambda R(w)]$ (штрафная форма) $\Leftrightarrow$ $\min \RSS(w)$ при $R(w) \le \beta$ (условная форма). Больше $\lambda$ --- меньше $\beta$.
\end{propbox}

\qsub{Геометрическая картина}
Линии уровня $\RSS$ --- эллипсы вокруг $w_{\text{OLS}}$. Область ограничения: L2 --- круг ($\norm{w}_2 \le \beta$), L1 --- ромб ($\norm{w}_1 \le \beta$, вершины на осях). Решение --- точка касания наименьшего эллипса с границей области.

\begin{propbox}{Почему L1 разрежает, а L2 нет}
\textbf{L1 (ромб)}: острые углы НА ОСЯХ (точки $(\beta,0), (0,\beta)$, где один вес $= 0$). Углы <<ловят>> эллипс широким диапазоном направлений подхода $\Rightarrow$ решение часто в углу $\Rightarrow$ зануление. \textbf{L2 (круг)}: гладкая граница, касание в точке, которая с вероятностью 1 не на оси $\Rightarrow$ оба веса ненулевые, просто сжаты. То же явление, что soft-threshold vs умножение (билет 14) --- геометрически.
\end{propbox}

\begin{fishkabox}
Ridge --- когда все признаки полезны и есть коллинеарность; Lasso --- когда нужен отбор из многих мусорных. Геометрически при коллинеарных ОБОИХ полезных признаках: круг не имеет привилегированных направлений $\to$ делит вес пропорционально, устойчиво к шуму; ромб навязывает дискретный выбор <<один признак = 0, другой = всё>>, причём выбор решает шум. L1 насильно вводит асимметрию туда, где данные симметричны.
\end{fishkabox}

\bilet{16}{Elastic Net и функция потерь Хубера}

\begin{defbox}{Elastic Net}
\[
  L = \tfrac1n\norm{Xw - y}^2 + \lambda_1\norm{w}_1 + \lambda_2\norm{w}_2^2.
\]
L1 даёт зануление, L2 --- устойчивость при коллинеарности. \textbf{Grouping effect}: $|w_i - w_j| \le \frac{\text{const}}{\lambda_2}\sqrt{2(1-\rho_{ij})}$ --- сильно коррелирующие признаки ($\rho_{ij} \to 1$) получают близкие веса, L2 <<склеивает>> группу, L1 зануляет её целиком. <<Lasso берёт один случайный из группы, Elastic Net --- всю группу>>.
\end{defbox}

\begin{defbox}{Функция потерь Хубера}
\[
  L_\delta(r) = \begin{cases} \tfrac12 r^2, & |r| \le \delta\\ \delta(|r| - \delta/2), & |r| > \delta \end{cases}, \quad r = y - \wh y.
\]
Малые остатки --- как MSE (гладко, хороший градиент), большие --- как MAE (выброс не доминирует). Непрерывна и дифференцируема в стыке $|r| = \delta$ (в отличие от чистой MAE).
\end{defbox}

\begin{propbox}{Предельные случаи Хубера}
$\delta \to 0$: почти вся область --- линейная ветка $\Rightarrow$ пропорциональна MAE. $\delta \to \infty$: условие $|r| \le \delta$ для всех $r$ $\Rightarrow$ только $\tfrac12 r^2$ = MSE. Чувствительность к выбросам: MSE $\gg$ Huber $>$ MAE.
\end{propbox}

\begin{fishkabox}
Почему не всегда Elastic Net / Huber: они добавляют гиперпараметры ($\lambda_1,\lambda_2$ / $\delta$) --- дороже подбор. MSE остаётся дефолтом: MLE при гауссовом шуме (ЦПТ), гладкость, простота. Elastic Net нужен при известной структуре коррелирующих групп; Huber --- при известных тяжёлых выбросах.
\end{fishkabox}

\bilet{17}{Выпуклость, L-гладкость и сильная выпуклость}

\begin{defbox}{Три свойства}
\begin{itemize}
  \item \textbf{Выпуклость}: $f(\lambda x + (1-\lambda)y) \le \lambda f(x) + (1-\lambda)f(y)$. Следствие: любой локальный минимум глобален.
  \item \textbf{L-гладкость}: $\norm{\nabla f(x) - \nabla f(y)} \le L\norm{x - y}$ (липшицев градиент). Для дважды диф.: $L = \lambda_{\max}(\nabla^2 f)$. Даёт безопасный шаг $h \le 1/L$.
  \item \textbf{$\mu$-сильная выпуклость}: $f(x) - \tfrac\mu2\norm{x}^2$ выпукла. Функция крутая везде, нет плато, минимум единственный.
\end{itemize}
\end{defbox}

\begin{rembox}{L-гладкость $\ne$ выпуклость}
L-гладкость --- ТОЛЬКО про скорость изменения градиента, НЕ гарантирует глобальность минимума. Функция потерь нейросети L-гладкая, но невыпуклая (много локальных минимумов). Гарантию <<локальный = глобальный>> дают только выпуклость и сильная выпуклость.
\end{rembox}

\qsub{МНК как пример}
МНК: выпукл, L-гладок с $L = \lambda_{\max}(X\tr X)$, но не обязательно сильно выпукл ($\mu = \lambda_{\min} \approx 0$ при плохой обусловленности). МНК+L2 (Ridge): $\mu$-сильно выпукл с $\mu = \tau$ (гессиан $2(X\tr X + \tau E)$, мин.\ собств.\ число $\ge 2\tau$).

\begin{fishkabox}
Зачем регуляризация, если она <<уводит от истинного $w$>>? МНК-оценка несмещена лишь В СРЕДНЕМ по многим выборкам, но у одной реальной выборки при мультиколлинеарности разброс (variance) огромен. Ridge жертвует малым bias ради резкого падения variance: аналогия --- стрелок B (систематически чуть смещён, но кучно) на одном выстреле ближе к цели, чем стрелок A (в среднем точен, но разброс огромный). Ошибка $=$ Bias$^2$ + Variance; окупается, когда variance доминирует.
\end{fishkabox}

\bilet{18}{Скорость сходимости градиентного спуска}

\begin{thmbox}{Геометрическая сходимость}
Для L-гладких $\mu$-сильно выпуклых функций GD (шаг $h \le 1/L$) сходится геометрически:
\[
  \norm{w^{(k+1)} - w^*}^2 \le q^k \norm{w^{(0)} - w^*}^2,\ q < 1, \qquad K \approx O\!\left(\frac{L}{\mu}\ln\frac1\varepsilon\right).
\]
\end{thmbox}

\qsub{Разбор формулы}
$\ln(1/\varepsilon)$ --- логарифмическая зависимость от точности: улучшить в 10 раз $\to$ не в 10 раз больше итераций, а константная добавка. $L/\mu$ --- число обусловленности задачи; для МНК $L/\mu = \lambda_{\max}/\lambda_{\min} = \cond(X\tr X)$. Большое $L/\mu$ --- вытянутая чаша (билет 5) --- зигзаг, много итераций. $L/\mu \approx 1$ --- сфера, быстро.

\qsub{Почему не за один шаг}
GD использует только градиент (1-й порядок) --- <<знает лишь наклон под ногами>>. Метод Ньютона (2-й порядок, гессиан) для квадратичной функции доходит за один шаг.

\begin{fishkabox}
Рост $\tau$ в Ridge: $L/\mu = (\lambda_{\max}+\tau)/\tau = \lambda_{\max}/\tau + 1$ строго УБЫВАЕТ ($d/d\tau < 0$). При $\tau \to \infty$: $L/\mu \to 1$ (идеальная обусловленность, быстрая сходимость). Но GD сходится к ДРУГОЙ, более смещённой точке (не $w_{\text{OLS}}$, а сжатое Ridge-решение). Скорость оптимизации и качество решения --- разные оси: регуляризация упрощает саму оптимизацию И меняет то, что ищем.
\end{fishkabox}

\bilet{19}{Стохастический градиентный спуск (SGD) и mini-batch}

\begin{defbox}{SGD}
Градиент по одному объекту / мини-батчу вместо всей выборки: $\nabla L_i(w) = 2 x_i(x_i\tr w - y_i)$. Итерация дёшева, но оценка градиента шумная.
\end{defbox}

\begin{propbox}{Несмещённость и дисперсия}
При случайном равновероятном выборе $i$: $\E[\nabla L_i(w)] = \nabla L(w)$ (несмещена), но большая дисперсия --- конкретный $\nabla L_i$ может сильно отличаться от истинного направления. Траектория SGD дрожит вокруг направления полного GD.
\end{propbox}

\begin{thmbox}{Условия Роббинса--Монро}
Для сходимости несмотря на шум расписание шага $\alpha_k$ должно удовлетворять
\[
  \sum_k \alpha_k = \infty \quad(\text{дотянуться до минимума}), \qquad \sum_k \alpha_k^2 < \infty \quad(\text{успокоиться, не дрожать вечно}).
\]
Пример: $\alpha_k = 1/k$ (гармонический ряд расходится, $\sum 1/k^2$ сходится). Только $\alpha_k = \text{const}$: вечное дрожание. Только $\alpha_k = 1/k^2$: сумма шагов конечна, до далёкого минимума не дойти.
\end{thmbox}

\qsub{Размер батча}
Малый: шумно, дёшево, шум помогает выпрыгивать и регуляризует. Большой: стабильно, дорого. Mini-batch (32--512): компромисс + параллелизм GPU.

\begin{fishkabox}
Два РАЗНЫХ эффекта одного шума, работают одновременно (не последовательно): (1) \textbf{выпрыгивание из плохой ямы} --- масштаб навигации по невыпуклому ландшафту (в какой долине окажемся); (2) \textbf{implicit regularization} --- шум не даёт дойти до острого/переобученного дна ВНУТРИ хорошей ямы, склоняет к широким flat minima (лучше обобщают). Первое --- <<в какой долине>>, второе --- <<насколько остро вниз в этой долине сядем>>.
\end{fishkabox}

\bilet{20}{Ускоренные и адаптивные методы: Momentum, Adam}

\begin{defbox}{Momentum (<<тяжёлый шарик>>)}
\[
  v = \gamma v + \eta g, \qquad w = w - v, \qquad g = \nabla f(w),\ \gamma \approx 0.9.
\]
$v$ --- накопленная скорость (память о прошлых шагах). Вдоль пологого направления (стабильный $g$) скорость накапливается --- разгон; вдоль крутого (знакопеременный $g$) вклады $+\eta g$ и $-\eta g$ гасят друг друга --- демпфирование зигзага. \textbf{NAG}: градиент в упреждающей точке $w - \gamma v$.
\end{defbox}

\begin{defbox}{Адаптивные методы (свой шаг каждому весу)}
\textbf{AdaGrad}: $S = S + g^2$, $w = w - \eta g/(\sqrt{S} + \varepsilon)$ --- шаг падает у весов с историей больших $g$; минус: $S$ растёт вечно, шаг затухает до нуля. \textbf{RMSProp}: $S = \beta S + (1-\beta)g^2$ --- EMA с забыванием, шаг не затухает навсегда.
\end{defbox}

\begin{defbox}{Adam = Momentum + RMSProp}
\begin{align*}
  m &= \beta_1 m + (1-\beta_1)g \quad(\text{направление}), & v &= \beta_2 v + (1-\beta_2)g^2 \quad(\text{масштаб}),\\
  \wh m &= \frac{m}{1 - \beta_1^t}, & \wh v &= \frac{v}{1 - \beta_2^t}, \qquad w = w - \eta\frac{\wh m}{\sqrt{\wh v} + \varepsilon}.
\end{align*}
$\beta_1 \approx 0.9$, $\beta_2 \approx 0.999$, $t$ --- номер итерации.
\end{defbox}

\begin{proofbox}{Зачем коррекция смещения}
$m_0 = 0 \Rightarrow m_1 = \beta_1\cdot 0 + (1-\beta_1)g_1 = (1-\beta_1)g_1$. При $\beta_1 = 0.9$: $m_1 = 0.1 g_1$ --- всего 10\% истинного градиента (сильно занижено). Коррекция: $\wh m_1 = m_1/(1 - \beta_1^1) = (1-\beta_1)g_1/(1-\beta_1) = g_1$ --- точно восстанавливает. При росте $t$ множитель $(1-\beta_1^t) \to 1$, коррекция автоматически выключается.
\end{proofbox}

\qsub{AdamW}
В обычном Adam L2-штраф $2\tau w$ тоже делится на $\sqrt{\wh v}$ $\to$ у весов с большим градиентом регуляризация ослабляется непредсказуемо. AdamW применяет weight decay $w \leftarrow w(1-\lambda)$ ОТДЕЛЬНО от адаптивной нормировки --- корректная L2.

\begin{fishkabox}
Adam/AdamW --- дефолт (быстро, мало настройки, устойчив к разномасштабным градиентам --- NLP, трансформеры). SGD+Momentum часто даёт ЛУЧШЕЕ финальное обобщение (CV) --- Adam иногда садится в острый минимум, а шум SGD ведёт к flat minima. State-of-the-art в CV нередко всё ещё SGD+momentum.
\end{fishkabox}

\bilet{21}{Разложение на смещение, дисперсию и шум (bias-variance)}

\begin{thmbox}{Bias--variance разложение}
Для $y = f(x) + \varepsilon$, $\E\varepsilon = 0$, $\Var\varepsilon = \sigma^2$, $\bar f(x) = \E_D[\wh f(x)]$:
\[
  \text{EPE}(x) = \underbrace{(f(x) - \bar f(x))^2}_{\text{Bias}^2} + \underbrace{\E_D[(\bar f(x) - \wh f(x))^2]}_{\text{Variance}} + \underbrace{\sigma^2}_{\text{Noise}}.
\]
\end{thmbox}

\begin{proofbox}{Обнуление перекрёстных членов}
Раскрываем $\E[((f - \bar f) + (\bar f - \wh f) + \varepsilon)^2]$. Перекрёстные члены зануляются: $(f - \bar f)$ --- константа; $\E[\bar f - \wh f] = 0$ по определению $\bar f = \E_D[\wh f]$; $\E\varepsilon = 0$ и $\varepsilon$ независим от $\wh f$. Остаётся сумма трёх квадратов.
\end{proofbox}

\qsub{Смысл слагаемых}
Bias$^2$ --- систематическая неспособность модели выразить $f$ (недообучение). Variance --- чувствительность к конкретной выборке (переобучение; связь с билетом 11 --- неустойчивость = высокая variance). Noise --- неустранимый предел ($\sigma^2$).

\begin{propbox}{Как зависят от сложности и $n$}
Рост сложности: Bias$\downarrow$, Variance$\uparrow$. Регуляризация ($\tau\uparrow$): Bias$\uparrow$, Variance$\downarrow$. \textbf{Bias} --- вопрос К МОДЕЛИ (форма/сложность), НЕ лечится ростом $n$. \textbf{Variance} --- вопрос К ВЫБОРКЕ, ПАДАЕТ с ростом $n$ (больше данных $\to$ выборки представительнее $\to$ $\wh f$ меньше прыгает, ЗБЧ).
\end{propbox}

\begin{fishkabox}
Learning curves: train-ошибка РАСТЁТ с $n$ (сложнее подогнать больше точек, $\to$ Bias$^2$+Noise сверху), val-ошибка ПАДАЕТ (Variance$\downarrow$, $\to$ тот же предел снизу), кривые сходятся к Bias$^2$+Noise. Разрыв между ними = мера Variance. Частая ошибка: сказать <<variance растёт с $n$>> --- наоборот, падает. Bias лечат сложнее моделью/признаками, Variance --- данными/регуляризацией/бэггингом. Бэггинг и рост $n$ снижают Variance почти бесплатно (не трогая Bias).
\end{fishkabox}

\bilet{22}{Энтропия, кросс-энтропия и KL-дивергенция}

\begin{defbox}{Три величины}
\[
  H(p) = -\sum_i p_i \log p_i, \quad H(p,q) = -\sum_i p_i \log q_i, \quad KL(p\Vert q) = \sum_i p_i\log\frac{p_i}{q_i} = H(p,q) - H(p).
\]
\end{defbox}

\qsub{Информационный смысл}
$H(p)$ --- среднее число бит для оптимального кодирования $x \sim p$ (выводится однозначно из аксиом непрерывности, монотонности, аддитивности --- теорема Шеннона). Максимум у равномерного (полная неопределённость). Честная монета: $H = 1$ бит. $H(p,q)$ --- реальное число бит, если кодируем оптимально под $q$ (модель), а данные приходят из $p$ (истина); $H(p,q) \ge H(p)$, равенство $\Leftrightarrow q = p$. $KL$ --- лишние биты из-за $q \ne p$; $KL \ge 0$, $= 0 \Leftrightarrow p = q$, НЕсимметрична (не метрика).

\qsub{Связи с ML}
Максимизация правдоподобия $\Leftrightarrow$ минимизация $KL(p_{\text{data}}\Vert p_{\text{model}})$. Кросс-энтропия = log-loss логрег/softmax. Энтропия = критерий ветвления деревьев (information gain).

\begin{exbox}{Асимметрия KL на 3 точках}
Истина $p = (0.5,\, 0,\, 0.5)$ (моды A, C; B невозможна). Кандидаты: $q_{\text{cover}} = (0.4, 0.2, 0.4)$ (покрывает обе моды, чуть <<мусорит>> в B), $q_{\text{single}} = (0, 0, 1)$ (игнорирует моду A).\\
$KL(p\Vert q_{\text{single}})$: точка A даёт $0.5\log(0.5/0) = \infty$. $KL(p\Vert q_{\text{cover}}) \approx 0.22$ (в B: $0\cdot\log(0/0.2) = 0$, штрафа нет). \textbf{Forward KL} наказывает игнор моды, не наказывает лишнюю массу.\\
$KL(q_{\text{single}}\Vert p) \approx 0.69$ (в A: $0\cdot\log = 0$, штрафа за игнор нет). $KL(q_{\text{cover}}\Vert p)$: в B $0.2\log(0.2/0) = \infty$. \textbf{Reverse KL} наказывает массу там, где истины нет.
\end{exbox}

\begin{fishkabox}
Ключ к асимметрии --- ВЕС перед логарифмом. $KL(p\Vert q) = \sum p_i\log(p_i/q_i)$ взвешен по $p$: штраф $\infty$ если $q \to 0$ там где $p > 0$; но $0$ если $p = 0$ (лишняя масса $q$ бесплатна) $\Rightarrow$ \textbf{mean-seeking}, накрывает все моды. $KL(q\Vert p)$ взвешен по $q$: штраф $\infty$ если $p \to 0$ там где $q > 0$; но $0$ если $q = 0$ (игнор моды $p$ бесплатен) $\Rightarrow$ \textbf{mode-seeking}, садится на одну моду.
\end{fishkabox}

\bilet{23}{Кросс-валидация и утечки данных}

\begin{defbox}{K-fold CV}
Разбить на $k$ фолдов; $k$ раз обучать на $k{-}1$, валидировать на оставшемся; усреднить. Каждый объект ровно раз в роли validation. Надёжнее одного split (усредняет по разбиениям).
\end{defbox}

\qsub{Разновидности}
\textbf{Stratified}: сохраняет доли классов (критично при дисбалансе). \textbf{LOO} ($k = n$): почти несмещённая оценка, но дорого ($n$ обучений) и высокая дисперсия (каждый фолд = один объект). \textbf{Time series split}: только прошлое $\to$ будущее.

\qsub{Nested CV}
Внешний цикл --- честная оценка, внутренний --- подбор гиперпараметров. Иначе (подбор на тех же фолдах, где меряем качество) --- оптимистичное смещение.

\begin{defbox}{Виды утечки данных (data leakage)}
\begin{enumerate}
  \item \textbf{Target leakage} --- признак есть следствие таргета (<<дата выписки>> при предсказании выписки).
  \item \textbf{Препроцессинг} --- статистики ($\mu,\sigma$, target encoding, импутация) считать ТОЛЬКО на train внутри фолда.
  \item \textbf{Временная} --- будущее $\to$ прошлое.
  \item \textbf{Дубликаты} train/test.
\end{enumerate}
Симптом --- подозрительно высокое качество, не подтверждающееся в проде.
\end{defbox}

\begin{fishkabox}
Тонкая утечка через препроцессинг в CV: если $\mu,\sigma$ посчитаны по ВСЕЙ выборке до разбиения, то стандартизируя train-часть фолда, мы используем $\mu,\sigma$, в которые внесли вклад объекты validation ЭТОГО фолда. Train-признаки статистически коррелируют со значениями validation $\Rightarrow$ модель косвенно <<подглядывает>>. Правильно: пересчитывать $\mu,\sigma$ заново для train-части КАЖДОГО фолда (не один раз глобально).
\end{fishkabox}

% =====================================================================
\razdel{3}{Оптимизация второго порядка}
% =====================================================================

\emph{Раздел для полноты; на собесах по classic ML спрашивают редко (кроме интуиции GD vs Ньютон vs BFGS). Изложено компактно.}

\bilet{24}{Метод Ньютона}
\[
  x_{k+1} = x_k - [\nabla^2 L(x_k)]^{-1}\nabla L(x_k).
\]
Квадратичная аппроксимация по Тейлору (учитывает кривизну $\nabla^2 L$ = гессиан). Для квадратичной функции (МНК) --- точный минимум за 1 шаг: <<видит форму горы целиком>>.

\bilet{25}{Проблемы Ньютона и Левенберг--Марквардт}
Минусы Ньютона: гессиан дорог ($O(m^3)$ на обращение), может быть необратим/плохо обусловлен, вдали от минимума расходится. \textbf{Левенберг--Марквардт} (нелинейный МНК): $(J\tr J + \lambda I)p = -J\tr f$. $\lambda I$ как Ridge гарантирует обратимость: большой $\lambda \to$ GD (надёжно, медленно), малый $\to$ Гаусс--Ньютон (быстро). $\lambda$ адаптируется по успеху шага.

\bilet{26}{Демпфированный Ньютон и линейный поиск}
$x_{k+1} = x_k + \gamma_k H^{-1}\nabla f$ --- размер шага $\gamma_k$ и приближённый (замороженный) $H$. Line search: при известном направлении $p_k$ задача сводится к одномерной $\gamma_k^* = \arg\min_\gamma f(x_k + \gamma p_k)$; неточный поиск --- условия Армихо/Вульфа (достаточное убывание).

\bilet{27}{Квазиньютоновские методы: секущее уравнение}
Строим приближение $B_k \approx [\nabla^2 f]^{-1}$ по истории градиентов (только 1-е производные). $s_k = x_{k+1} - x_k$, $y_k = \nabla f(x_{k+1}) - \nabla f(x_k)$: секущее уравнение $s_k = B_{k+1}y_k$ (недоопределено $\to$ доп.\ требования: симметричность, $\succ 0$, близость к $B_k$).

\bilet{28}{Метод SR-1 (Symmetric Rank-1)}
Добавка ранга 1: $B_{k+1} = B_k + \dfrac{(s_k - B_k y_k)(s_k - B_k y_k)\tr}{(s_k - B_k y_k)\tr y_k}$. При $k \to \infty$ сходится к обратному гессиану. Минус: может терять положительную определённость $\to$ направление не убывания.

\bilet{29}{Алгоритм BFGS и L-BFGS}
Добавка ранга 2. Гарантия: если $B_0 \succ 0$ и $y_k\tr s_k > 0$ (обеспечивает line search) --- все $B_k \succ 0$, направление всегда убывающее. \textbf{L-BFGS}: хранит не всю матрицу, а последние $\sim$10--20 пар $(s_k, y_k)$ --- для больших размерностей.

\begin{fishkabox}
Иерархия GD $\to$ BFGS $\to$ Ньютон = сколько информации о кривизне (0/1/2-й порядок): GD --- только наклон (дёшево, много шагов), Ньютон --- точная кривизна (дорого, мало шагов), BFGS --- приближает кривизну по истории градиентов (средне). Почему не всегда BFGS: (1) $O(m^2)$ память --- невозможно для нейросетей с млрд параметров; (2) плохо совместим с mini-batch --- шумная история градиентов ломает оценку кривизны. Поэтому в deep learning --- SGD/Adam, а BFGS/L-BFGS --- в классическом ML (логрег, CRF) с full-batch.
\end{fishkabox}

% =====================================================================
\razdel{4}{Снижение размерности}
% =====================================================================

\bilet{30}{Теорема о наилучшем низкоранговом приближении (Эккарта--Янга)}

\begin{thmbox}{Эккарт--Янг}
\[
  \min_{B\colon \rank(B) = k} \norm{X - B} \quad\Longrightarrow\quad B = \sum_{j=1}^k u_j\sigma_j v_j\tr, \qquad \norm{X - B} = \sigma_{k+1}.
\]
Наилучшее приближение ранга $k$ (по норме Фробениуса/спектральной) = усечённое SVD (top-$k$ троек). Ошибка = первое отброшенное сингулярное число.
\end{thmbox}

\qsub{Зачем снижать размерность}
Борьба с шумом (малые $\sigma$ = шумовые направления); визуализация; скорость; борьба с мультиколлинеарностью; проклятие размерности (kNN деградирует).

\begin{fishkabox}
Оптимальность ГЛОБАЛЬНАЯ, не поточечная. $\norm{X - B}_F^2 = \sum_i\norm{x_i - b_i}^2$ --- минимизируется СУММА по объектам, не каждое слагаемое. Нетипичный объект (вариация в отброшенных направлениях с малым $\sigma$) реконструируется плохо, даже если $B$ оптимальна в среднем. Пример: PCA на лицах ловит общие вариации (освещение, поворот), но редкая особенность (шрам) уходит в отброшенные компоненты. Связь с выбросами (билет 94): низкоранговое приближение <<стирает>> редкие сигналы.
\end{fishkabox}

\bilet{31}{PCA через максимизацию дисперсии}

\begin{defbox}{Задача PCA}
Найти ортогональное направление $\text{Rot}$ ($\norm{\text{Rot}} = 1$) с максимальной дисперсией проекций: $\max_{\text{Rot}}\norm{X\cdot\text{Rot}}^2$ при $\norm{\text{Rot}} = 1$.
\end{defbox}

\begin{proofbox}{Вывод через Лагранжиан}
$\mathcal{L} = \text{Rot}\tr(X\tr X)\text{Rot} - \lambda(\text{Rot}\tr\text{Rot} - 1)$. Градиент:
\[
  \nabla\mathcal{L} = 2(X\tr X)\text{Rot} - 2\lambda\text{Rot} = 0 \;\Longrightarrow\; (X\tr X)\text{Rot} = \lambda\text{Rot}.
\]
Направления максимальной дисперсии = собственные векторы $X\tr X$ = столбцы $V$ (SVD), $\lambda = \sigma^2$ = дисперсии проекций.
\end{proofbox}

\qsub{Эквивалентность двух взглядов}
<<Максимум дисперсии вдоль $\text{Rot}$>> $\Leftrightarrow$ <<минимум ошибки реконструкции $\perp \text{Rot}$>> (билет 30). Полная энергия $\sum\norm{x_i}^2 = $ const раскладывается по теореме Пифагора на объяснённую дисперсию + ошибку реконструкции; максимизация первого минимизирует второе.

\begin{fishkabox}
PCA (unsupervised, максимум дисперсии) vs LDA (supervised, максимум разделимости классов). Направление максимальной дисперсии может быть ПЕРПЕНДИКУЛЯРНО направлению разделения классов: два длинных узких облака классов, вытянутых вдоль одной оси, но сдвинутых перпендикулярно. PCA выберет ось вытянутости (максимум дисперсии), спроецирует на неё --- классы схлопнутся и перемешаются. LDA нашёл бы ось разделения. Мораль: PCA-препроцессинг перед классификацией может УНИЧТОЖИТЬ разделимость.
\end{fishkabox}

\vspace{1cm}
\begin{center}\small\color{gray}\itshape Конспект продолжается по мере прохождения билетов (Разделы 5--13).\end{center}

\end{document}
